% Configuration PDF commune aux documents Husson.
\usepackage{fontspec}
\usepackage{xcolor}
\usepackage{fvextra}
\usepackage{newunicodechar}
\usepackage{enumitem}
\usepackage{etoolbox}
\usepackage{needspace}
\usepackage{pdflscape}

\definecolor{sectionblue}{HTML}{1D4E89}
\definecolor{commentgray}{HTML}{6F6F6F}
\definecolor{pythonaccent}{HTML}{9A6100}
\definecolor{pythonfill}{HTML}{FFF5D9}
\definecolor{raccent}{HTML}{276DC3}
\definecolor{rfill}{HTML}{EAF2FB}

\newunicodechar{≥}{\ensuremath{\geq}}
\newunicodechar{≤}{\ensuremath{\leq}}
\setlength{\emergencystretch}{3em}

\makeatletter
\renewcommand*{\@pnumwidth}{2.8em}
\renewcommand*{\@tocrmarg}{3.6em}
\makeatother

\setkomafont{disposition}{\rmfamily}
\setkomafont{section}{\rmfamily\color{sectionblue}\bfseries}
\RedeclareSectionCommand[
  beforeskip=0.8\baselineskip,
  afterskip=0.35\baselineskip
]{section}

\setlength{\parindent}{0pt}
\setlength{\parskip}{3pt}
\setlist[itemize]{itemsep=1pt,parsep=1pt,topsep=2pt}
\setlist[enumerate]{itemsep=1pt,parsep=1pt,topsep=2pt}

% Ajoute 0,4 cm au-dessus et au-dessous du contenu de chaque ligne de tableau.
% La boîte typographique normale est rétablie avant de composer les paragraphes
% multiligne : la marge n'est donc appliquée qu'une fois par cellule.
\renewcommand{\arraystretch}{1}
\setlength{\extrarowheight}{0pt}
\makeatletter
\newbox\hussontablestrutbox
\newcommand{\hussontablepadding}{%
  \setbox\hussontablestrutbox=\copy\strutbox
  \setbox\strutbox=\hbox{%
    \vrule width 0pt
      height \dimexpr\ht\hussontablestrutbox+0.4cm\relax
      depth \dimexpr\dp\hussontablestrutbox+0.4cm\relax
  }%
  \let\hussontablesavedmkpream\@mkpream
  \def\@mkpream##1{%
    \hussontablesavedmkpream{##1}%
    \setbox\strutbox=\copy\hussontablestrutbox
    \let\@mkpream\hussontablesavedmkpream
  }%
}
\makeatother

% Les tableaux gardent leur hauteur naturelle, avec une taille homogène.
\AtBeginEnvironment{longtable}{\small\hussontablepadding}
\AtBeginEnvironment{tabular}{\small\hussontablepadding}

% Les blocs de code sont compacts et peuvent revenir à la ligne.
\DefineVerbatimEnvironment{Highlighting}{Verbatim}{%
  fontsize=\footnotesize,
  breaklines,
  breakanywhere,
  commandchars=\\\{\}
}
\providecommand{\CommentTok}[1]{\textcolor{commentgray}{\textit{#1}}}

\newcommand{\codelanguagelabel}[2]{%
  \Needspace{5\baselineskip}%
  \par\addvspace{0.55\baselineskip}%
  \noindent\begingroup
  \setlength{\fboxsep}{3pt}%
  \colorbox{#1}{\color{white}\sffamily\bfseries\scriptsize #2}%
  \endgroup\par\nobreak\vspace{-0.25\baselineskip}%
}

\newenvironment{pythoncode}{%
  \codelanguagelabel{pythonaccent}{PYTHON}%
  \begingroup\colorlet{shadecolor}{pythonfill}%
}{%
  \endgroup\par\addvspace{0.25\baselineskip}%
}

\newenvironment{rcode}{%
  \codelanguagelabel{raccent}{R}%
  \begingroup\colorlet{shadecolor}{rfill}%
}{%
  \endgroup\par\addvspace{0.25\baselineskip}%
}

\makeatletter
\renewcommand{\maketitle}{
  \begin{center}
    {\Large\bfseries\rmfamily \@title\par}
    \vskip 0.35em
    {\large\rmfamily \@subtitle\par}
    \vskip 0.65em
    {\normalsize\rmfamily \@author\par}
  \end{center}
}
\makeatother

\renewcommand{\contentsname}{}
\AtBeginDocument{
  \let\oldtableofcontents\tableofcontents
  \renewcommand{\tableofcontents}{
    \begingroup
    \footnotesize
    \setlength{\parskip}{1pt}
    \oldtableofcontents
    \endgroup
  }
}
